\loai{Một vật nhận nhiệt lượng từ bên ngoài}
\begin{vidu}
	Người ta thực hiện thí nghiệm xác định nhiệt dung riêng của đồng với một miếng đồng có khối lượng 850 g. Lúc đầu nhiệt độ của miếng đồng là $12^\circ C$. Thời gian từ khi bật độ phận đốt nóng có công suất 40 W đến khi nhiệt độ của miếng đồng tăng tới $30^\circ C$ là 190 giây. Xác định nhiệt dung riêng của đồng.
	\loigiai{
		\begin{itemize}
			\item 
			\item 
			\item 
			\item 
			\item 
			\item 
		\end{itemize}
	}
\end{vidu}
\loai{Hai vật trao đổi nhiệt - Tính nhiệt độ}
\begin{vidu}%[1P6-1.2T][Loai3]
	Để xác định nhiệt độ của một cái lò, người ta đưa vào một miếng sắt có khối lượng $m=22,3$ g. Khi miếng sắt có nhiệt độ bằng nhiệt độ của lò, người ta lấy ra và thả ngay vào nhiệt lượng kế chứa 450 g nước ở 15$^{\circ}C$, nhiệt độ của nước tăng lên tới 22,5$^{\circ}C$. Biết $c_{Fe}=478$ J/(kg.K), $c_{H_2O}=4180$ J/(kg.K). Xác định nhiệt độ của lò, bỏ qua sự hấp thụ nhiệt của nhiệt lượng kế.
	\choice
	{$1870,9^{\circ}C$}
	{\True $1345,9^{\circ}C$}
	{$2640,9^{\circ}C$}
	{$7370,6^{\circ}C$}
	\loigiai{ 
		\Binhluan{
			\begin{bang}[tabularx=X|X|X]{}
				&Chưa trao đổi&Cân bằng nhiệt\\\hline
				Nước&$t_1=15^\circ C$&$t=22,5^\circ C$\\\hline
				Sắt&$t_2=?$&$t=22,5^\circ C$
			\end{bang}
		}
		{$Q=mc.(t_s-t_{tr})$.\\
			$Q>0$ vật thu nhiệt.\\
			$Q<0$ vật toả nhiệt.
		}
		\begin{itemize} 
			\item Nhiệt lượng trao đổi của sắt: $Q_{Fe}=m_{Fe}.c_{Fe}(t-t_2)=239,8-10,7t_2$ J.
			\item Nhiệt lượng trao đổi của nước: $Q_{H_2O}=m_{H_2O}.c_{H_2O}(t-t_1)=14107,5$ J.
			\item Theo điều kiện cân bằng nhiệt: 
			\begin{align*}
				Q_{Fe}+Q_{H_2O}=0 \Rightarrow m_{Fe}.c_{Fe}(t_2-t)+239,8-10,7t_2=0\Rightarrow t_2=1345,9^{\circ}C.
			\end{align*}
		\end{itemize}
	}
\end{vidu}
\loai{Hai vật trao đổi nhiệt - Tính nhiệt dung riêng}
\begin{vidu}
	Một khay sắt có khối lượng 1,2 kg được cách điện và làm nóng bằng máy sưởi 500 W trong 4 phút. Nhiệt độ của khay tăng từ $22^\circ C$ đến $45^\circ C$. Xác định nhiệt dung riêng của sắt. Bỏ qua mất mát nhiệt lượng do môi trường.
	\loigiai{
		\begin{itemize}
			\item 
			\item 
			\item 
			\item 
			\item 
			\item 
		\end{itemize}
	}
\end{vidu}

\loai{Ba vật trao đổi nhiệt - Tính nhiệt độ khi cân bằng}
\begin{vidu}%[1P6-1.2T][Loai6] 
	Một cốc nhôm có khối lượng 120 g chứa 400 g nước ở nhiệt độ $24\doC$. Người ta thả vào cốc nước một thìa đồng khối lượng 80 g ở nhiệt độ $100\doC$. Xác định nhiệt độ của nước trong cốc khi có sự cân bằng nhiệt.  Bỏ qua sự trao đổi nhiệt với môi trường. Biết nhiệt dung riêng của nhôm là 880 J/(kg.K), của nước là $4,19.10^3$ J/(kg.K), của đồng là 380 J/(kg.K).
	\choice
	{\True $25,27\doC$}
	{$14,5\doC$}
	{$24,67\doC$}
	{$26,73\doC$}
	\loigiai{
		\begin{itemize} 
			\item Nhiệt lượng trao đổi của cốc nhôm: $Q_{Al} = m_{Al}c_{Al}(t - t_1)$.
			\item Nhiệt lượng trao đổi của nước: $Q_{H_2O} = m_{H_2O}c_{H_2O}(t - t_1)$.
			\item Nhiệt lượng trao đổi của thìa đồng: $Q_{Cu} = m_{Cu}c_{Cu}(t - t_2)$.
			\item Khi có cân bằng nhiệt: 
			$Q_{Al} +Q_{H_2O}+Q_{Cu}=0 \Rightarrow t = 25,27^{\circ}C.$
	\end{itemize}}
\end{vidu}
\loai{Ba vật trao đổi nhiệt - Tính khối lượng}
\begin{vidu}%[1P6-1.2T][Loai1]
	Một nhiệt lượng kế bằng nhôm chứa nước có khối lượng tổng cộng $0,5$ kg ở $t_1 = 25^{\circ}$C. Cho vào nhiệt lượng kế một quả cân bằng đồng khối lượng $m_2 = 0,4$ kg ở $t_2 = 100^{\circ}$C. Khi có cân bằng nhiệt, nhiệt độ của nước $t = 35^{\circ}$C. Bỏ qua sự truyền nhiệt ra môi trường bên ngoài. Tìm khối lượng của nhiệt kế và của nước. Nhiệt dung riêng của nhôm $c = 880$ J/(kg.K), của nước $c_1 = 4200$ J/(kg.K), của đồng $c_2 = 380$ J/(kg.K).
	\choice
	{$335\ g;\ 265\ g$}
	{\True $335\ g;\ 165\ g$}
	{$435\ g;\ 165\ g$}
	{$435\ g;\ 265\ g$}
	\loigiai{
		\Binhluan{
			\begin{itemize} 
				\item Gọi x là khối lượng của nhiệt lượng kế, khối lượng của nước là: $0,5 - x$.
				\item Nhiệt lượng do nhiệt lượng kế và nước trao đổi: 
				\begin{align*}
					Q_1 = [x.c + (0,5 - x)c_1] (t - t_1) = 100 (210 - 332x).
				\end{align*}
				\item Nhiệt lượng do quả cân trao đổi: $ Q_2 = m_2c_2 (t - t_2) = -9880$ J.
				\item Khi có cân bằng nhiệt: 
				\begin{align*}
					Q_1 + Q_2=0 \Rightarrow 100(210 - 332x) = 9880 \Rightarrow x \approx 0,335 \ kg.
				\end{align*}
				\item Khối lượng nhiệt lượng kế: $335\ g$, của nước: $165\ g$.
			\end{itemize}
		}
		{Áp dụng điều kiện cân bằng nhiệt và dùng SHIFT SOLVE cho ra kết quả nhanh chóng}
	}
\end{vidu}
\loai{Ba vật trao đổi nhiệt - Tính nhiệt dung riêng}
\begin{vidu}%[1P6-1.2T][Loai8] 
	Một nhiệt lượng kế bằng đồng khối lượng $m_1 = 100$ g chứa $m_2 = 375$ g nước ở nhiệt độ $25\doC$. Cho vào nhiệt lượng kế một vật bằng kim loại khối lượng $m_3 = 400$ g ở $90\doC$. Biết nhiệt độ khi có sự cân bằng nhiệt là $30\doC$. Bỏ qua sự trao đổi nhiệt với môi trường. Tìm nhiệt dung riêng của miếng kim loại. Cho biết nhiệt dung riêng của đồng là 380 J/(kg.K) và của nước là 4200 J/(kg.K).
	\choice
	{633 J/kg. K}
	{\True 336 J/kg. K}
	{362 J/kg. K}
	{880 J/kg. K}
	\loigiai{
		\Binhluan{
			\begin{itemize} 
				\item Nhiệt lượng trao đổi của nhiệt lượng kế: $Q_{NLK} = m_{Cu}c_{Cu}(t - t_1)$.
				\item Nhiệt lượng trao đổi của nước: $Q_{H_2O} = m_{H_2O}c_{H_2O}(t - t_1)$.
				\item Nhiệt lượng trao đổi của kim loại: $Q_{Kl} = m_{Kl}c_{Kl}(t - t_2)$.
				\item Khi có cân bằng nhiệt: 
				$Q_{NLK} +Q_{H_2O}+Q_{Kl}=0 \Rightarrow c_{Kl} = 336\ J/(kg.K).$
			\end{itemize}
		}
		{Đặt ẩn số cần tìm là x và dùng SHIFT SOLVE}
	}
\end{vidu}

\loai{Hai vật trao đổi nhiệt - Có nhiệt lượng bên ngoài cấp vào}
\begin{vidu}%[1P6-1.2G][Loai10]
	Một ấm đun nước bằng nhôm có có khối lượng 400 g, chứa 3 lít nước được đun trên bếp. Khi nhận được nhiệt lượng 740 kJ thì ấm đạt đến nhiệt độ $80^{\circ}C$.  Bỏ qua sự trao đổi nhiệt với môi trường bên ngoài. Nhiệt độ ban đầu của ấm là bao nhiêu? Biết $c_{A1}=880$ J/(kg.K), $c_{H_2O}=4190$ J/(kg.K).
	\choice
	{8,15$^{\circ}C$}
	{8,15 K}
	{\True 22,7$^{\circ}C$}
	{22,7 K}
	\loigiai{ 
		\Binhluan{
			\begin{itemize}
				\item Nhiệt lượng trao đổi của nước: 			
				\begin{align*}
					Q_{H_2O}=m_{H_2O}.c_{H_2O}(t-t_1)=1005600-12570t_1.
				\end{align*}
				\item Nhiệt lượng trao đôi của ấm nhôm: 
				\begin{align*}
					Q_{Al}=m_{Al}.c_{Al}(t-t_1)=28160-352t_1.
				\end{align*}
				\item Vì ấm và nước nhận được nhiệt lượng 740 kJ, ta có:
				\begin{align*}
					Q_{H_2O}+Q_{Al}=740.10^3\Rightarrow t_1=22,7^{\circ}C.
				\end{align*}
			\end{itemize}
		}
		{Đây không phải là quá trình cân bằng nhiệt}
	}
\end{vidu}
\loai{Thực hiện công}
\begin{vidu}
	Một người cọ xát một miếng sắt dẹt có khối lượng 150 g trên một sàn nhà bê tông. Sau một khoảng thời gian, miếng sắt nóng thêm $12^\circ C$. Tính công mà người này đã thực hiện, giả sử rằng 40\% công đó được dùng để làm nóng miếng sắt. Biết nhiệt dung riêng của sắt là 460 J/kg.K.
	\loigiai{
		\begin{itemize}
			\item 
			\item 
			\item 
			\item 
			\item 
			\item 
		\end{itemize}
	}
\end{vidu}
\loai{Áp dụng thêm nguyên lí I}
\begin{vidu}
	Một vật khối lượng 1 kg trượt không vận tốc đầu từ đỉnh xuống chân một mặt phẳng nghiêng dài 21 m, nghiêng $30^\circ$ so với phương ngang. Tốc độ ở chân mặt phẳng là 4,1 m/s. Tính công của lực ma sát và độ biến thiên nội năng của vật trong quá trình chuyển động trên. Bỏ qua sự trao đổi nhiệt của vật với mặt phẳng nghiêng. Lấy $g=9,8\ m/s^2$.
	\loigiai{
		\begin{itemize}
			\item 
			\item 
			\item 
			\item 
			\item 
			\item 
		\end{itemize}
	}
\end{vidu}
\loai{Tổng hợp}

\begin{vidu}
	Vận động viên chạy Marathon mất rất nhiều nước trong khi thi đấu. Các vận động viên thường chỉ có thể chuyển hoá khoảng 20\% năng lượng hoá học dự trữ trong cơ thể thành năng lượng dùng cho các hoạt động cơ thể, đặc biệt là hoạt động chạy. Phần năng lượng còn lại chuyển thành nhiệt thải ra ngoài nhờ sự bay hơi của nước qua hô hấp và da để giữ cho nhiệt độ của cơ thể không đổi. Nếu vận động viên dùng hết 11000 kJ trong cuộc thi thì có khoảng bao nhiêu lít nước đã thoát ra khỏi cơ thể? Coi nhiệt độ cơ thể của vận động viên hoàn toàn không đổi và nhiệt hoá hơi riêng của nước ở nhiệt độ của vận động viên là $2.10^6\ J/kg$.	
	\loigiai{
		\begin{itemize}
			\item 
			\item 
			\item 
			\item 
			\item 
			\item 
		\end{itemize}
	}
\end{vidu}
